\documentclass[11pt]{article}
\usepackage{airxiv}
\usepackage{amsmath}
\usepackage{array}
\usepackage{booktabs}
\usepackage{float}
\usepackage{enumitem}

\graphicspath{{../figures/}}

% Machine-generated statistics. Regenerate with
%   us-gdp-regimes export-report-numbers --config config/default.yaml
% Every quoted number below is a macro defined in this file, so the prose can
% never drift from the model outputs in data/models/.
\input{generated_numbers.tex}

\title{
  Growth Regimes in the United States since 1920:\\
  Structural Breaks, Fiscal Structure, and the Productivity--Pay Divergence
}

\author{
  Diogo Ribeiro\thanks{ESMAD, Instituto Polit\'{e}cnico do Porto. All estimates,
  tables, and figures are reproducible from the accompanying open-source pipeline
  (\texttt{us-gdp-regime-1920}); every quantitative statement in the text is
  generated programmatically from the model outputs. Correspondence:
  \texttt{dfr@esmad.ipp.pt}.}\\
  \small ESMAD, Instituto Polit\'{e}cnico do Porto
}

\date{July 10, 2026}

\begin{document}
\maketitle

\begin{abstract}
We characterise the annual growth process of United States real GDP from 1920
onward as a sequence of statistically identified regimes rather than a single
trend. Using a Maddison-based long-run real GDP series validated against BEA
national accounts, we first show that log GDP is empirically indistinguishable
from a unit-root process while its growth rate is stationary, so growth---not the
level---is the appropriate object for regime analysis. A dynamic-programming
segmentation with information-criterion selection identifies $\NumGlobalRegimes$
global regimes; because a single pooled variance lets the extreme 1929--1945
episode dominate, we apply a recursive refinement that re-segments each regime on
its own scale, yielding $\NumRegimes$ regimes and resolving the long postwar era
into a high-growth $\PostwarEarlyStart$--$\PostwarEarlyEnd$ phase and the
post-$\PostwarSplitYear$ slowdown. Parametric-bootstrap break tests, break-date
confidence intervals, and a robustness grid over tuning choices establish which
breaks are firm and which are fragile. We then situate the regimes against
federal fiscal structure, the dynamic response of growth to narrative tax-regime
shocks, and the widening gap between real GDP per capita and real median
earnings. The exercise is deliberately descriptive: it treats statistical breaks
as structure to be interpreted, not as identified causal effects.
\end{abstract}

\keywords{United States GDP, growth regimes, structural breaks, unit roots, local
projections, fiscal policy, real wages, reproducible research}

\vspace{0.3em}
\noindent\textbf{JEL classification:} C22, E32, N12, O47, E62.

\reportnote{Reproducible from the repository pipeline; raw downloaded data are
excluded from version control. Every statistic in the text is a machine-generated
macro derived from the model outputs in \texttt{data/models/}.}

\section{Introduction}

The long-run growth of the United States economy is routinely summarised by a
single average rate. That summary is analytically convenient but substantively
misleading: a century that contains the Great Depression, wartime mobilisation,
the postwar boom, and the post-2000 slowdown is not well described by one number.
This paper asks a deliberately simple question---did the United States grow at
one roughly constant speed after 1920, or is its growth history better read as a
sequence of distinct regimes?---and answers it with a transparent, fully
reproducible empirical workflow.

Our approach has four components. First, we establish the correct object of
study. A linear trend fits log real GDP almost perfectly in sample
($R^2 = \TrendRSquared$), but that fit is a well-known artefact of persistence:
augmented Dickey--Fuller \citep{saiddickey1984} and KPSS \citep{kpss1992} tests
together indicate that the log level is consistent with a unit root while the
growth rate is stationary. Regimes are therefore estimated on growth. Second, we
detect regimes with an exact dynamic-programming segmentation of the growth
series and select their number by information criterion. Because a single global
variance lets the violent 1929--1945 years dominate, we introduce a recursive
refinement that re-segments each regime on its own variance scale; this is what
separates the postwar boom from the post-$\PostwarSplitYear$ slowdown that a
pooled model conceals. Third, we quantify the uncertainty of the regimes with
sequential structural-break tests in the spirit of Bai and Perron
\citep{baiperron1998,baiperron2003}, bootstrap confidence intervals for the break
dates, and a robustness grid across tuning choices. Fourth, we relate the regimes
to federal fiscal structure, to the dynamic response of growth to narrative
tax-regime shocks estimated by local projections \citep{jorda2005}, and to the
divergence between aggregate output per capita and the earnings of the median
worker.

\paragraph{Contributions.} The paper makes three modest contributions. (i) It
documents that the standard deterministic-trend summary of U.S.\ growth is
statistically inappropriate, and provides the unit-root and stationarity evidence
that motivates a growth-regime treatment. (ii) It shows that a naive global
segmentation understates postwar heterogeneity because of variance pooling, and
that a simple recursive refinement---corroborated by three independent criteria
and by a subsample break test---recovers the post-2000 slowdown as a distinct
regime. (iii) It integrates fiscal, tax-dynamic, and distributional evidence into
the same reproducible framework while maintaining a strict separation between
descriptive timing, lagged association, and causal identification.

\paragraph{Related literature.} The empirical strategy draws on the multiple
structural-change framework of \citet{baiperron1998,baiperron2003} and on the
observation of \citet{perron1989} that unmodelled breaks bias unit-root tests
toward non-rejection---a caveat we take seriously in interpreting the level
diagnostics. Dynamic responses are estimated by local projections
\citep{jorda2005} with Newey--West covariance \citep{neweywest1987}; the
distinction between plausibly exogenous and endogenous tax changes follows the
narrative-identification tradition of \citet{romerromer2010}. Substantively, the
post-2000 regime connects to the debate over the productivity slowdown
\citep{gordon2016,fernald2015} and secular stagnation \citep{summers2015}, while
the output--earnings divergence connects to work on the productivity--pay gap
\citep{stansburysummers2018} and rising earnings inequality \citep{autor2014}.
The long-run series is the Maddison Project Database \citep{boltvanzanden2020}.

\paragraph{Roadmap.} Section~\ref{sec:data} describes the data and validation.
Section~\ref{sec:methods} sets out the empirical strategy. Section~\ref{sec:results}
presents the trend diagnostics, regimes, break inference, and postwar
decomposition; Section~\ref{sec:robustness} reports robustness. Sections~\ref{sec:fiscal}
and \ref{sec:distribution} add the fiscal, tax-dynamic, and distributional layers.
Section~\ref{sec:discussion} discusses the findings, Section~\ref{sec:limitations}
the limitations, and Section~\ref{sec:conclusion} concludes.

\section{Data}
\label{sec:data}

\subsection{Core GDP series and validation}

The primary series is annual United States real GDP from 1920 onward, constructed
from the Maddison Project Database 2023 \citep{boltvanzanden2020} as real GDP per
capita multiplied by population. Maddison is designed for long-run historical
comparison and is the natural source for a 1920 starting point; the growth-rate
analysis is invariant to the constant unit rescaling implied by the population
units. The modern benchmark is the Bureau of Economic Analysis annual real GDP
series (FRED \texttt{GDPCA}), available from 1929.

Over the $\OverlapStart$--$\OverlapEnd$ overlap ($\OverlapYears$ annual
observations), Maddison-derived and BEA growth rates correlate at
$\GrowthCorrelation$, with a mean absolute difference of $\GrowthMad$ and a root
mean squared difference of $\GrowthRmse$ percentage points. The largest
discrepancies (Appendix~\ref{app:validation}) fall in the wartime years, when
historical and national-accounts output concepts diverge most. The agreement is
close enough for regime-level analysis while reminding the reader that the two
sources are distinct objects.

\subsection{Fiscal, tax, and distributional data}

Three further layers provide context. The fiscal layer uses FRED/OMB federal
debt, receipts, outlays, deficit, and interest-outlay ratios as a share of GDP.
The tax layer uses a curated catalog of broad federal tax-law changes with signed
ordinal shock values and narrative classifications distinguishing plausibly
exogenous long-run reforms from deficit-driven, wartime, and countercyclical
changes. The distributional layer adds BEA real GDP per capita, BLS real median
weekly earnings and real hourly compensation, the composition of federal
receipts, statutory top and bottom income-tax rates, and BLS Consumer Expenditure
federal income-tax rates by income quintile. These series answer different
questions and are not interchangeable welfare measures.

\section{Empirical Strategy}
\label{sec:methods}

\subsection{Trend and unit-root diagnostics}

The long-run trend is the ordinary least squares regression of log real GDP on a
linear time trend,
\begin{equation}
  \log(GDP_t) = \alpha + \beta\, t + \epsilon_t,
  \label{eq:trend}
\end{equation}
with the annualised trend rate $e^{\beta}-1$. Because annual macroeconomic
residuals are serially correlated, $\beta$ is reported with Newey--West
heteroskedasticity- and autocorrelation-consistent standard errors
\citep{neweywest1987}. A high $R^2$ in \eqref{eq:trend} is not evidence for a
deterministic trend: a near-integrated series is mechanically well fit by a line.
We therefore test the log level and the growth rate with the augmented
Dickey--Fuller test \citep{saiddickey1984}, whose null is a unit root, and the
KPSS test \citep{kpss1992}, whose null is stationarity. We read the two jointly,
and we note the caution of \citet{perron1989} that unmodelled structural breaks
bias ADF toward failing to reject a unit root.

\subsection{Growth regimes and recursive refinement}

Regimes are a piecewise-constant mean model for annual growth,
\begin{equation}
  g_t = \mu_j + u_t, \qquad t \in \text{regime } j,
  \label{eq:regime}
\end{equation}
where the segmentation into contiguous regimes is chosen to minimise a Gaussian
information criterion. For a partition into $k$ segments with within-segment sum
of squared residuals $\mathrm{SSR}_k$ and $n$ observations, the criterion is
\begin{equation}
  \mathrm{IC}(k) = n \log\!\left(\frac{\mathrm{SSR}_k}{n}\right) + p_k \, c(n),
  \qquad p_k = 2k,
  \label{eq:ic}
\end{equation}
with $c(n) = \log n$ for BIC and $c(n) = 2$ for AIC; the parameter count $p_k=2k$
comprises $k$ segment means, $k-1$ break locations, and one residual variance. The
globally optimal partition for each $k$ is obtained exactly by dynamic
programming over prefix sums, and $k$ is chosen to minimise \eqref{eq:ic}. Each
segment is labelled ``above'' or ``below'' the full-sample mean; the labels are
descriptive.

A single global segmentation imposes one residual variance on the whole century.
This is a poor assumption for U.S.\ growth: the 1929--1945 swings inflate the
pooled variance so much that calmer postwar mean shifts never clear the penalty.
We therefore apply a \emph{recursive refinement}. Starting from the global fit,
each segment is re-segmented on its own observations---and thus its own residual
variance---and split whenever \eqref{eq:ic} supports it, up to a bounded depth.
This is the headline model; the global fit is retained as a reference.

\subsection{Break inference}

The number of regimes is tested, not assumed. For each step from $l$ to $l+1$
segments we form a sequential $\mathrm{supF}(l+1\mid l)$ statistic from the
optimal $\mathrm{SSR}$ of the two nested models. Because the break date is
estimated, the statistic does not follow a standard $F$ distribution; we obtain
its p-value by a parametric bootstrap ($\BreakTestBootstrap$ replications) that
simulates Gaussian data from the fitted smaller model \citep{baiperron1998}.
Uncertainty in the break \emph{locations} is quantified separately by a residual
bootstrap ($\BreakCiBootstrap$ replications) that re-segments each resample at a
fixed number of breaks and reports percentile intervals for each break year.

\subsection{Tax-regime dynamics and distributional indices}

Dynamic responses of growth to tax-regime shocks are estimated by local
projections \citep{jorda2005},
\begin{equation}
  g_{t+h} = \alpha_h + \beta_h\, s_t + \gamma_h\, g_{t-1}
  + \gamma'_h\, g_{t-2} + e_{t+h}, \qquad h = 0,\dots,10,
  \label{eq:lp}
\end{equation}
where $s_t$ is a signed tax-regime shock. Because the horizon-$h$ residual is an
MA($h$) process by construction, the Newey--West lag length grows with $h$. We
estimate \eqref{eq:lp} on all catalogued events and, separately, on the subset
classified as plausibly exogenous \citep{romerromer2010}. Distributional series
are indexed to a common base year, $\mathrm{index}_{x,t} = 100\, x_t / x_{base}$,
and the headline gap is $\mathrm{index}_{GDPpc,t} - \mathrm{index}_{earnings,t}$.

\section{Results}
\label{sec:results}

\subsection{Trend and unit-root diagnostics}

Table~\ref{tab:trend} reports the log-trend regression. The slope is precisely
estimated even under HAC covariance ($t = \TrendHacT$), implying annualised trend
growth of $\TrendGrowth\%$. As anticipated, this precision reflects persistence
rather than a genuine deterministic trend. Table~\ref{tab:unitroot} makes the
point formal: for the log level, ADF does not reject a unit root
($\LevelAdfStat$, $p \LevelAdfP$) and KPSS rejects trend-stationarity
($p \LevelKpssP$); for growth, the verdicts reverse (ADF $p \GrowthAdfP$; KPSS
$p \GrowthKpssP$). Growth is the stationary object, and the appropriate one for
regime analysis.

\begin{table}[H]
  \centering
  \caption{Log-trend regression of real GDP, 1920--present}
  \label{tab:trend}
  \begin{tabular}{lr}
    \toprule
    \input{generated_trend_reg_table.tex}
  \end{tabular}
  \\[0.4em]
  {\footnotesize\begin{minipage}{0.86\linewidth}
  \textit{Notes.} OLS of $\log$ real GDP on a linear time trend,
  Eq.~\eqref{eq:trend}. Newey--West standard error in parentheses. The high $R^2$
  is an artefact of persistence, not evidence of a deterministic trend.
  \end{minipage}}
\end{table}

\begin{table}[H]
  \centering
  \caption{Unit-root and stationarity tests}
  \label{tab:unitroot}
  \begin{tabular}{lccccc}
    \toprule
    Series & Spec. & ADF stat & ADF $p$ & KPSS stat & KPSS $p$ \\
    \midrule
    \input{generated_unit_root_table.tex}
  \end{tabular}
  \\[0.4em]
  {\footnotesize\begin{minipage}{0.86\linewidth}
  \textit{Notes.} ADF null: unit root; KPSS null: stationarity. Specification
  \texttt{ct} includes a constant and trend; \texttt{c} a constant only. Read
  jointly, the level is difference-stationary and growth is stationary.
  \end{minipage}}
\end{table}

Figure~\ref{fig:trend} plots log real GDP with the fitted trend of
Eq.~\eqref{eq:trend}. The line tracks the level closely, but the diagnostics
above show that this apparent fit is the signature of a near-integrated series
rather than evidence of a stable deterministic path.

\begin{figure}[H]
  \centering
  \includegraphics[width=\linewidth]{log_real_gdp_trend.png}
  \caption{Log real GDP and fitted linear trend. The close fit reflects
  persistence, not a deterministic trend (cf.\ Tables~\ref{tab:trend} and
  \ref{tab:unitroot}). Source: Maddison Project Database 2023.}
  \label{fig:trend}
\end{figure}

\subsection{Growth regimes}

Table~\ref{tab:regimes} reports the $\NumRegimes$ regimes of the headline model,
plotted in Figure~\ref{fig:regimes}. The 1929--1933 contraction
($\approx -5.8\%$) and the 1934--1944 rebound and mobilisation ($\approx +8.1\%$)
dominate the early sample. A plain global fit leaves the entire post-1950 era as a
single block ($\NumGlobalRegimes$ regimes); the recursive refinement splits it
into a high-growth $\PostwarEarlyStart$--$\PostwarEarlyEnd$ regime averaging
$\PostwarEarlyMean\%$ and a slower $\PostwarLateStart$--$\PostwarLateEnd$ regime
averaging $\PostwarLateMean\%$, against a full-sample mean of
$\LongRunMeanGrowth\%$.

\begin{table}[H]
  \centering
  \caption{Estimated annual growth regimes (headline recursive model)}
  \label{tab:regimes}
  \begin{tabular}{rrrrl}
    \toprule
    Start & End & Years & Mean growth & Label \\
    \midrule
    \input{generated_regime_table.tex}
  \end{tabular}
\end{table}

\begin{figure}[H]
  \centering
  \includegraphics[width=\linewidth]{gdp_growth_regimes.png}
  \caption{Annual real GDP growth with piecewise-constant regimes. Bands are
  statistical segments, not claims that any single event caused an entire segment.
  Source: Maddison Project Database 2023.}
  \label{fig:regimes}
\end{figure}

\subsection{Break significance and dating}

Table~\ref{tab:breaktests} reports the sequential break tests for the global
model. Additional regimes earn small bootstrap p-values up to
$\NumGlobalRegimes$ segments and are insignificant thereafter. This is the formal
counterpart of the global information-criterion choice, and it explains why a
single-variance model stops before splitting the postwar era: relative to the
enormous 1929--1945 residual variance, the postwar mean shift does not clear the
global bar. That is a property of the pooled-variance assumption, not evidence
that the postwar era is homogeneous.

\begin{table}[H]
  \centering
  \caption{Sequential $\mathrm{supF}(l+1\mid l)$ break tests}
  \label{tab:breaktests}
  \begin{tabular}{rrrr}
    \toprule
    Segments (null) & Segments (alt.) & $F$ statistic & Bootstrap $p$ \\
    \midrule
    \input{generated_break_test_table.tex}
  \end{tabular}
  \\[0.4em]
  {\footnotesize\begin{minipage}{0.72\linewidth}
  \textit{Notes.} Parametric-bootstrap p-values, $\BreakTestBootstrap$
  replications, simulating from the fitted smaller model.
  \end{minipage}}
\end{table}

The break \emph{dates} are unevenly identified (Table~\ref{tab:breakcis}). The
Depression and mobilisation breaks are sharp, whereas the start of the long
postwar regime near $\WidestBreakYear$ has a percentile interval spanning
$\WidestBreakCiLow$--$\WidestBreakCiHigh$. That imprecision is itself a symptom:
the global model is forcing a heterogeneous postwar era through one break and one
mean.

\begin{table}[H]
  \centering
  \caption{Bootstrap confidence intervals for the break years (global model)}
  \label{tab:breakcis}
  \begin{tabular}{rrcr}
    \toprule
    Break & Point year & 90\% interval & Std.\ (years) \\
    \midrule
    \input{generated_break_ci_table.tex}
  \end{tabular}
\end{table}

\subsection{Postwar decomposition}

Re-segmenting the postwar era on its own variance scale resolves the ambiguity.
Table~\ref{tab:postwar} shows that three independent choices---the postwar
subsample under BIC, the postwar subsample under AIC, and the full sample under
AIC---all split the era at $\PostwarSplitYear$ into the same
$\PostwarEarlyStart$--$\PostwarEarlyEnd$ and $\PostwarLateStart$--$\PostwarLateEnd$
regimes. On the postwar subsample the single break at $\PostwarSplitYear$ is
statistically significant ($p \PostwarBreakP$) while no further postwar break is.
This is the post-2000 growth slowdown, and it is the regime that recursive
refinement adds to Table~\ref{tab:regimes}.

\begin{table}[H]
  \centering
  \caption{Postwar decomposition: three criteria agree on the
  $\PostwarSplitYear$ split}
  \label{tab:postwar}
  \begin{tabular}{lcccc}
    \toprule
    Method & Early regime & Mean & Late regime & Mean \\
    \midrule
    \input{generated_postwar_table.tex}
  \end{tabular}
\end{table}

\section{Robustness}
\label{sec:robustness}

Table~\ref{tab:robustness} re-estimates the break years across
$\RobustnessNScenarios$ scenarios: minimum segment sizes of 4, 5, 7, and 10
years; the BIC and AIC criteria; and exclusion of the 1941--1945 wartime window.
Break years are clustered within a two-year tolerance and reported with the
number of scenarios in which they appear. The Depression trough
($\MostRobustBreakYear$) recurs in all $\MostRobustBreakScenarios$ scenarios, and
the late-1920s, mid-1940s, and late-1940s breaks recur in the large majority. The
post-2000 slowdown appears only under the AIC scenario in this global grid,
consistent with the finding above that a pooled-variance model under-selects it;
its support comes instead from the subsample decomposition of
Table~\ref{tab:postwar}. The message is twofold: the pre-war regime structure is
highly robust to tuning, while the postwar split is real but criterion-dependent
and is best identified on its own scale.

\begin{table}[H]
  \centering
  \caption{Recurring break years across robustness scenarios}
  \label{tab:robustness}
  \begin{tabular}{ccc}
    \toprule
    Representative year & Range & Scenarios \\
    \midrule
    \input{generated_robustness_table.tex}
  \end{tabular}
  \\[0.4em]
  {\footnotesize\begin{minipage}{0.78\linewidth}
  \textit{Notes.} Scenarios vary the minimum segment size (4, 5, 7, 10), the
  criterion (BIC, AIC), and exclude 1941--1945. Break years within two years are
  clustered.
  \end{minipage}}
\end{table}

\section{Fiscal Structure and Tax-Regime Dynamics}
\label{sec:fiscal}

\subsection{Fiscal context}

Federal fiscal ratios move sharply across wars, recessions, and crises, but their
same-year comovement with growth is modest: gross debt-to-GDP correlates with
growth at $\GrossDebtCorrSame$ contemporaneously and $\GrossDebtCorrLagOne$ at a
one-year lag, with smaller correlations for receipts, outlays, and deficits
(Figure~\ref{fig:fiscal}). These are timing devices, not identified fiscal
effects: debt, deficits, and receipts are endogenous to the same shocks that move
output.

\begin{figure}[H]
  \centering
  \includegraphics[width=\linewidth]{fiscal_context.png}
  \caption{Federal debt, budget ratios, and annual GDP growth. A timing and
  context device; it does not identify fiscal causality.}
  \label{fig:fiscal}
\end{figure}

\subsection{Dynamic response to tax-regime shocks}

Table~\ref{tab:lp} and Figure~\ref{fig:taxlp} report local-projection estimates
of growth following signed tax-regime shocks, for all catalogued events and for
the plausibly exogenous subset \citep{romerromer2010}. The coefficients are
small, switch sign across horizons, and are almost never individually
significant; confidence intervals are wide, especially for the exogenous subset,
which contains few events. The results should be read as a dynamic screening
exercise, not as tax multipliers. Their value is methodological: by allowing
delayed responses, they show that the data do not support a simple same-year tax
story, consistent with anticipation, phase-ins, and heterogeneous channels.

\begin{table}[H]
  \centering
  \caption{Local projections of GDP growth on tax-regime shocks}
  \label{tab:lp}
  \begin{tabular}{rcccc}
    \toprule
    & \multicolumn{2}{c}{All events} & \multicolumn{2}{c}{Exogenous subset} \\
    \cmidrule(lr){2-3}\cmidrule(lr){4-5}
    Horizon $h$ & Coef. & (SE) & Coef. & (SE) \\
    \midrule
    \input{generated_lp_table.tex}
  \end{tabular}
  \\[0.4em]
  {\footnotesize\begin{minipage}{0.82\linewidth}
  \textit{Notes.} Coefficient $\beta_h$ from Eq.~\eqref{eq:lp}; Newey--West
  standard errors with lag growing in $h$. $^{*}$, $^{**}$, $^{***}$ denote
  significance at 10\%, 5\%, and 1\%.
  \end{minipage}}
\end{table}

\begin{figure}[H]
  \centering
  \includegraphics[width=\linewidth]{tax_local_projections.png}
  \caption{Local-projection estimates after signed tax-regime shocks. The
  restricted line uses only plausibly exogenous long-run reforms.}
  \label{fig:taxlp}
\end{figure}

\section{The Productivity--Pay Divergence}
\label{sec:distribution}

The regime that closes the sample---the post-2000 slowdown---coincides with a
long-running divergence between aggregate output and worker earnings.
Figure~\ref{fig:wagegap} indexes real GDP per capita, real median weekly
earnings, and real hourly compensation to $\IndexBaseYear$. By $\DistLatestYear$
real GDP per capita reached $\GdppcIndexLatest$ (more than double its base), while
real median weekly earnings reached only $\MedianEarningsIndexLatest$; real hourly
compensation ($\HourlyCompIndexLatest$) lies between. Aggregate output per person
can rise substantially while the median worker's measured real earnings grow
slowly.

\begin{figure}[H]
  \centering
  \includegraphics[width=\linewidth]{wage_gdp_gap.png}
  \caption{Real GDP per capita, real median weekly earnings, and real hourly
  compensation, indexed to $\IndexBaseYear$. The widening gap shows why GDP
  growth is not equivalent to broad worker buying power.}
  \label{fig:wagegap}
\end{figure}

The composition of the tax system shifted over the same period
(Figure~\ref{fig:taxburden}). Social-insurance contributions rose from
$\SocialInsuranceShareFirst\%$ of the selected federal receipt base in
$\ReceiptFirstYear$ to $\SocialInsuranceShareLatest\%$ in $\ReceiptLatestYear$,
while corporate income taxes fell from $\CorporateShareFirst\%$ to
$\CorporateShareLatest\%$. Federal income-tax progressivity by quintile, however,
did not collapse: the top quintile's average federal income-tax rate exceeded the
bottom-four-quintile average by $\QuintileSpreadFirst$ points in
$\QuintileFirstYear$ and $\QuintileSpreadLatest$ points in $\QuintileLatestYear$.
These proxies are not a full incidence model; they indicate a shift toward
social-insurance financing rather than a simple reduction in income-tax
progressivity.

\begin{figure}[H]
  \centering
  \includegraphics[width=\linewidth]{tax_burden_shift.png}
  \caption{Tax-burden shift proxies: receipt composition and, where available,
  average federal income-tax rates by income quintile.}
  \label{fig:taxburden}
\end{figure}

\section{Discussion}
\label{sec:discussion}

Three findings deserve emphasis. First, the appropriate summary of U.S.\ growth
is a stationary process punctuated by regimes, not a deterministic trend; the
unit-root evidence is not a technicality but the reason a regime treatment is
warranted, and it echoes the caution of \citet{perron1989} that breaks and unit
roots are easily confounded. Second, the modelling choice of one global variance
versus regime-specific variances is consequential: it is the difference between a
homogeneous 73-year postwar block and a postwar boom followed by a distinct
post-2000 slowdown. That slowdown is precisely the phenomenon at the centre of the
productivity-slowdown \citep{gordon2016,fernald2015} and secular-stagnation
\citep{summers2015} debates, and locating it as a data-driven break rather than an
assumed date is a useful discipline. Third, the divergence between output per
capita and median earnings shows that aggregate growth is an incomplete welfare
statistic; the pattern is consistent with the broader productivity--pay
\citep{stansburysummers2018} and inequality \citep{autor2014} literatures, though
our proxies cannot decompose its sources. Throughout, we resist causal readings:
the fiscal and tax layers are endogenous to the same shocks that drive output, and
the local-projection estimates are screening devices, not multipliers.

\section{Limitations}
\label{sec:limitations}

Three limitations bound the interpretation. \emph{Identification}: breakpoints,
fiscal ratios, tax laws, and wages all move with macroeconomic conditions; the
narrative classification of tax shocks mitigates but does not solve this, and the
event sample is small. \emph{Measurement}: Maddison is not identical to BEA
national accounts, median weekly earnings are not household income or total
compensation, and the tax-burden proxies are not a full incidence model.
\emph{Coverage}: the Maddison-based regime sample ends in $\PostwarEnd$, while
several FRED series in the fiscal and distributional layers extend to
$\DistLatestYear$ and may include a partial final year; coverage is therefore
reported by series rather than assumed uniform.

\section{Conclusion}
\label{sec:conclusion}

The United States did not grow at one constant speed after 1920. A deterministic
trend fits the log level in sample but is statistically inappropriate, because the
level is indistinguishable from a unit-root process while growth is stationary.
Read through growth, the century resolves into a small number of statistically
supported regimes: a 1920s expansion, the Depression, wartime mobilisation,
postwar demobilisation, the postwar boom, and the post-$\PostwarSplitYear$
slowdown. Recovering that last regime requires taking regime-specific variances
seriously; a pooled-variance model conceals it. Placed next to fiscal structure,
tax-regime dynamics, and the widening output--earnings gap, the regimes make a
simple point precisely: aggregate GDP growth is informative but incomplete, and
is best read as a trend-plus-regimes process reported alongside the fiscal and
distributional facts it omits.

\section*{Reproducibility Statement}

The paper is generated from the repository pipeline:
\begin{itemize}[leftmargin=1.5em]
  \item \texttt{us-gdp-regimes run --config config/default.yaml}
  \item \texttt{us-gdp-regimes make-fiscal-context --config config/default.yaml}
  \item \texttt{us-gdp-regimes make-tax-effects --config config/default.yaml}
  \item \texttt{us-gdp-regimes make-distributional-analysis --config config/default.yaml}
  \item \texttt{us-gdp-regimes export-report-numbers --config config/default.yaml}
\end{itemize}
The final command regenerates \texttt{generated\_numbers.tex} and the generated
tables that this paper inputs, so every quoted statistic is reconciled with the
model outputs. Raw downloaded data, model CSVs, figures, and generated tables are
reproducible artefacts and are not committed as source.

\appendix

\section{Source Validation Detail}
\label{app:validation}

Table~\ref{tab:sourcediff} lists the years with the largest absolute difference
between Maddison-derived and BEA (\texttt{GDPCA}) annual growth. The largest gaps
fall in the World War II years, where wartime output measurement and historical
reconstruction diverge most; excluding these years leaves close agreement across
the overlap.

\begin{table}[H]
  \centering
  \caption{Largest Maddison--BEA annual growth differences}
  \label{tab:sourcediff}
  \begin{tabular}{rrrr}
    \toprule
    Year & Maddison (\%) & BEA (\%) & Difference (pp) \\
    \midrule
    \input{generated_source_diff_table.tex}
  \end{tabular}
\end{table}

\begin{figure}[H]
  \centering
  \includegraphics[width=\linewidth]{fred_maddison_growth_comparison.png}
  \caption{Overlapping annual real GDP growth from Maddison-derived GDP and BEA
  \texttt{GDPCA}. A validation diagnostic, not a causal test.}
  \label{fig:sourcecompare}
\end{figure}

\begin{thebibliography}{99}

\bibitem[Autor(2014)]{autor2014}
Autor, D. H.
\newblock Skills, education, and the rise of earnings inequality among the
``other 99 percent''.
\newblock \emph{Science}, 344(6186):843--851, 2014.

\bibitem[Bai and Perron(1998)]{baiperron1998}
Bai, J. and Perron, P.
\newblock Estimating and testing linear models with multiple structural changes.
\newblock \emph{Econometrica}, 66(1):47--78, 1998.

\bibitem[Bai and Perron(2003)]{baiperron2003}
Bai, J. and Perron, P.
\newblock Computation and analysis of multiple structural change models.
\newblock \emph{Journal of Applied Econometrics}, 18(1):1--22, 2003.

\bibitem[Bolt and van Zanden(2020)]{boltvanzanden2020}
Bolt, J. and van Zanden, J. L.
\newblock Maddison style estimates of the evolution of the world economy: A new
2020 update.
\newblock Maddison Project Working Paper WP-15, University of Groningen, 2020.

\bibitem[Fernald(2015)]{fernald2015}
Fernald, J. G.
\newblock Productivity and potential output before, during, and after the Great
Recession.
\newblock \emph{NBER Macroeconomics Annual}, 29(1):1--51, 2015.

\bibitem[Gordon(2016)]{gordon2016}
Gordon, R. J.
\newblock \emph{The Rise and Fall of American Growth: The U.S. Standard of Living
Since the Civil War}.
\newblock Princeton University Press, 2016.

\bibitem[Jord\`a(2005)]{jorda2005}
Jord\`a, \`O.
\newblock Estimation and inference of impulse responses by local projections.
\newblock \emph{American Economic Review}, 95(1):161--182, 2005.

\bibitem[Kwiatkowski et al.(1992)]{kpss1992}
Kwiatkowski, D., Phillips, P. C. B., Schmidt, P., and Shin, Y.
\newblock Testing the null hypothesis of stationarity against the alternative of
a unit root.
\newblock \emph{Journal of Econometrics}, 54(1--3):159--178, 1992.

\bibitem[Newey and West(1987)]{neweywest1987}
Newey, W. K. and West, K. D.
\newblock A simple, positive semi-definite, heteroskedasticity and
autocorrelation consistent covariance matrix.
\newblock \emph{Econometrica}, 55(3):703--708, 1987.

\bibitem[Perron(1989)]{perron1989}
Perron, P.
\newblock The Great Crash, the oil price shock, and the unit root hypothesis.
\newblock \emph{Econometrica}, 57(6):1361--1401, 1989.

\bibitem[Romer and Romer(2010)]{romerromer2010}
Romer, C. D. and Romer, D. H.
\newblock The macroeconomic effects of tax changes: Estimates based on a new
measure of fiscal shocks.
\newblock \emph{American Economic Review}, 100(3):763--801, 2010.

\bibitem[Said and Dickey(1984)]{saiddickey1984}
Said, S. E. and Dickey, D. A.
\newblock Testing for unit roots in autoregressive-moving average models of
unknown order.
\newblock \emph{Biometrika}, 71(3):599--607, 1984.

\bibitem[Stansbury and Summers(2018)]{stansburysummers2018}
Stansbury, A. and Summers, L. H.
\newblock Productivity and pay: Is the link broken?
\newblock NBER Working Paper 24165, 2018.

\bibitem[Summers(2015)]{summers2015}
Summers, L. H.
\newblock Demand side secular stagnation.
\newblock \emph{American Economic Review: Papers \& Proceedings},
105(5):60--65, 2015.

\bibitem[Maddison Project Database(2023)]{maddison2023}
Maddison Project Database.
\newblock Maddison Project Database 2023.
\newblock \url{https://www.rug.nl/ggdc/historicaldevelopment/maddison/}.

\bibitem[FRED GDPCA]{fredgdpca}
Federal Reserve Bank of St. Louis.
\newblock Real Gross Domestic Product, annual, \texttt{GDPCA}.
\newblock \url{https://fred.stlouisfed.org/series/GDPCA}.

\bibitem[U.S. OMB]{ombtables}
Office of Management and Budget.
\newblock Historical Tables.
\newblock \url{https://www.whitehouse.gov/omb/budget/historical-tables/}.

\bibitem[BLS Consumer Expenditure Surveys]{blsce}
U.S. Bureau of Labor Statistics.
\newblock Consumer Expenditure Surveys.
\newblock \url{https://www.bls.gov/cex/}.

\bibitem[BEA National Accounts]{bea}
U.S. Bureau of Economic Analysis.
\newblock National Income and Product Accounts.
\newblock \url{https://www.bea.gov/data/gdp/gross-domestic-product}.

\end{thebibliography}

\end{document}
